\section*{第 \ref{ch:g9:electric-power-energy} 章 --- 电功率与能量}

\begin{solution}{exo:g9:electric-power-energy:1}
功率是能量转化的快慢——即每秒转化的焦耳数，
单位为瓦特：它像水龙头的流量，而不是水桶里的水。$P = U \times I$，
并有 $E = P \times t$。
\end{solution}

\begin{solution}{exo:g9:electric-power-energy:2}
灯泡：$230 \times 0.022 \approx \qty{5}{W}$。起动机：$12
\times 100 = \qty{1200}{W}$。
\end{solution}

\begin{solution}{exo:g9:electric-power-energy:3}
电热水壶：$2300 \div 230 = \qty{10}{A}$；微波炉：$700 \div
230 \approx \qty{3}{A}$。
\end{solution}

\begin{solution}{exo:g9:electric-power-energy:4}
因为供暖是能量供应中最昂贵的形式：要提高
水或房间的温度，所耗焦耳的速度，
是灯具和电子设备远远无法企及的，因此
每台取暖器的铭牌都以千瓦标示。
\end{solution}

\begin{solution}{exo:g9:electric-power-energy:5}
$\qty{1}{kWh} = 1000 \, \text{W} \times 3600 \, \text{s} =
\qty{3.6e6}{J}$。按焦耳计，一个月的用电量是十位数；
电表制造商选择了这样一个单位，使家庭
每月的用电量保持在几百的范围内。
\end{solution}

\begin{solution}{exo:g9:electric-power-energy:6}
取暖器：$2 \times 0.5 = \qty{1}{kWh}$。电视机：$0.1
\times 5 = \qty{0.5}{kWh}$。半小时的取暖器耗电
是整个晚上看电视用电量的两倍。
\end{solution}

\begin{solution}{exo:g9:electric-power-energy:7}
$E = 0.05 \times 4 \times 30 = \qty{6}{kWh}$；费用为 $6 \times
0.25 = 1.5$ 个货币单位——一个月的作业用电，
只需一块糕点的价钱。
\end{solution}

\begin{solution}{exo:g9:electric-power-energy:8}
电能表计量的是进入房屋的每一千瓦时。账单
用本月读数减去上月读数，
再乘以电价。
\end{solution}

\begin{solution}{exo:g9:electric-power-energy:9}
老式电器：$0.6 \times 4 \times 30 = \qty{72}{kWh}$。
现代电器：$0.05 \times 4 \times 30 = \qty{6}{kWh}$。节省：
\qty{66}{kWh}——每月节省 $16.5$ 个货币单位，
只需在梯子上忙活一个下午。
\end{solution}

\begin{solution}{exo:g9:electric-power-energy:10}
烘干机：$2000 \div 230 \approx \qty{8.7}{A}$。加上
电热水壶的 \qty{10}{A}，合计 $18.7$——超过了 \qty{16}{A}
支路的额定电流：断路器（理所当然地）会切断电路，
结束这次实验。
\end{solution}

\begin{solution}{exo:g9:electric-power-energy:11}
电热水壶：$60 \div 2.3 \approx 26$ 小时。灯泡：$60 \div
0.005 = 12000$ 小时——约合一年零四个月。然而
同一电池组只能驱动汽车行驶几百千米：
移动一吨重的机器所消耗的能量，
其规模之大，家庭照明相比之下微不足道。
\end{solution}

\begin{solution}{exo:g9:electric-power-energy:12}
预测：$t = E/P = \num{5.0e5} \div 2300 \approx
\qty{217}{s}$——约三分半钟。实际的
电热水壶用的时间更长：部分焦耳漏失在水之外，
——加热了壶身本身和周围空气，逸出的
早期水蒸气。能量从未消失，但也并非
全部送达预定之处：能量链条的漏失，
终于在你自家厨房里被测量到。
\end{solution}

\begin{solution}{pb:g9:electric-power-energy:1}
\textbf{1.} $2.3 \times 0.1 \times 30 = \qty{6.9}{kWh}$。
\textbf{2.} 电视机：$0.1 \times 4 \times 30 =
\qty{12}{kWh}$；灯具：$0.05 \times 5 \times 30 =
\qty{7.5}{kWh}$。
\textbf{3.} 洗衣机：$2 \times 1 \times 15 = \qty{30}{kWh}$；
热水器：$2.5 \times 2 \times 30 = \qty{150}{kWh}$。
\textbf{4.} $0.010 \times 720 = \qty{7.2}{kWh}$——隐藏的
十瓦功率，相当于一台电热水壶整月的用电量。
\textbf{5.} 平均功率 $120 \div 3 = \qty{40}{W}$；一个月：
$0.04 \times 720 = \qty{28.8}{kWh}$。
\textbf{6.} 总计 $\approx \qty{242}{kWh}$。主要耗电的是：
热水器（\qty{150}{}）、洗衣机（\qty{30}{}）、
冰箱（\qty{28.8}{kWh}）——加热类和常开类电器。
\textbf{7.} $242 \times 0.25 \approx 60.6$ 个货币单位；热水器
所占份额：$150 \times 0.25 = 37.5$——远超
电费的一半。
\textbf{8.} $2150 + 242 = 2392$ 个货币单位。
\textbf{9.} 新热水器线路：$2.5 \times 1.5 \times 30 =
\qty{112.5}{kWh}$；节省 \qty{37.5}{kWh}，约 $9.4$
个货币单位（每月）。
\textbf{10.} 一个闲置充电器：$0.0002 \times 720 =
\qty{0.14}{kWh}$——约百分之四个货币
单位；方案一每台充电器节省的是这个数的两百倍，
台充电器。审计能找到真正的节能杠杆，口号却只盯着充电器。
\textbf{11.} 电热水壶线路减半至 \qty{3.45}{kWh}：节省
\qty{3.45}{kWh}，不足一个货币单位——奶奶的
习惯是对的，而且幅度不大。
\textbf{12.} 例如：“千瓦时主要消耗在
水被加热的地方，以及机器从不休眠的地方。待机
耗电提醒我们时间的作用：微小的功率，
在 $t$ 始终不停时，也会推高电费。最划算的改变是缩短
淋浴时间——热水器主导着电费，因此它的使用时长
决定了节省量。”
\end{solution}
